En el artículo, se expone que actualmente transitamos por una cuarta revolución industrial.
Menciona que este periodo se caracteriza por una producción sin precedentes de información.
Destaca dos mecanismos principales de esta producción de información: por la vía del Internet de las Cosas (IoT), donde miles de millones de sensores y dispositivos interconectados recopilan detalles de nuestro entorno y nuestra actividad física.
Y por otro lado, la producción de datos por medio de las redes sociales, donde los individuos exponen voluntariamente sus ideas, experiencias y hábitos diariamente para interactuar en línea y evitar la exclusión social.

Toda esta hiperabundancia de datos ha desencadenado una profunda mercantilización de los mismos, transformando a los usuarios en activos altamente codiciados y redituables para las grandes compañías tecnológicas.

El texto detalla que la valuación de esta información se lleva a cabo de diversas formas, abarcando desde los miles de millones de dólares en ingresos publicitarios corporativos hasta ejercicios experimentales donde se subastan historiales de navegación individuales.
No obstante, existe una importante disparidad en la percepción y perspectiva de la gente sobre el valor económico de su Información Personalmente Identificable (PII); diversos estudios muestran que, aunque las personas otorgan mayor sensibilidad e importancia a datos como su ubicación o finanzas, la mayoría sigue cediendo de forma voluntaria sus hábitos y gustos en las redes sociales.

El texto subraya como conclusión la enorme necesidad de que la gente conozca cómo funciona este mercado y reflexione activamente sobre el valor monetario de sus datos.
La autora enfatiza que ceder nuestra privacidad no debe ser la única forma de pagar por el acceso a plataformas digitales, por lo que resulta indispensable que los individuos cuenten con alternativas tecnológicas que los empoderen.
De este modo, se busca que los usuarios tengan el control absoluto sobre su información para decidir cuándo y cómo compartirla, pero sobre todo, para protegerse cuando deseen que sus datos no se mercantilicen a favor de intereses ajenos.
